\documentclass[12pt,a4paper]{exam}
\usepackage[utf8]{inputenc}
\usepackage{geometry}
\geometry{margin=2cm}
\usepackage{enumitem}

\title{Questionário – Tema 4: Introdução à Meta-análise}
\author{Departamento de Física, Universidade Eduardo Mondlane}
\date{\today}

\begin{document}
\maketitle

\begin{questions}

%-----------------------------------
\question O que é meta-análise?  
\begin{choices}
\choice Uma revisão narrativa de estudos qualitativos
\CorrectChoice Uma técnica estatística que combina resultados de múltiplos estudos quantitativos
\choice Um teste de hipótese único aplicado a um estudo
\choice Um resumo de opiniões de especialistas
\end{choices}

%-----------------------------------
\question Qual é o objetivo principal de uma meta-análise?  
\begin{choices}
\choice Avaliar apenas um estudo isolado
\CorrectChoice Combinar evidências de vários estudos para aumentar o poder estatístico
\choice Produzir textos literários sobre um tema
\choice Criar figuras e gráficos sem análise estatística
\end{choices}

%-----------------------------------
\question Qual destas revisões não é quantitativa?  
\begin{choices}
\choice Revisão sistemática quantitativa
\choice Meta-análise
\CorrectChoice Revisão narrativa/literatura
\choice Revisão de métodos mistos
\end{choices}

%-----------------------------------
\question Qual das seguintes situações é mais apropriada para realizar uma meta-análise?  
\begin{choices}
\choice Quando há apenas um estudo disponível
\CorrectChoice Quando múltiplos estudos avaliam a mesma intervenção
\choice Quando se quer fazer uma análise de opinião
\choice Quando os estudos são completamente heterogêneos e incompatíveis
\end{choices}

%-----------------------------------
\question O que significa “tamanho do efeito” (effect size)?  
\begin{choices}
\choice Quantidade de participantes em um estudo
\CorrectChoice Medida padronizada da magnitude de um efeito observado entre grupos
\choice O número de gráficos gerados
\choice O nível de significância de um teste
\end{choices}

%-----------------------------------
\question Qual medida de efeito é usada para desfechos binários?  
\begin{choices}
\choice Diferença de médias (MD)
\choice Hazard ratio (HR)
\CorrectChoice Risco Relativo (RR) ou Odds Ratio (OR)
\choice Diferença de médias padronizada (SMD)
\end{choices}

%-----------------------------------
\question Qual medida de efeito é usada para desfechos contínuos?  
\begin{choices}
\choice Risco Relativo (RR)
\CorrectChoice Diferença de médias (MD) ou Diferença de médias padronizada (SMD)
\choice Odds Ratio (OR)
\choice Hazard ratio (HR)
\end{choices}

%-----------------------------------
\question O que um Forest Plot mostra?  
\begin{choices}
\choice A heterogeneidade dos estudos apenas
\CorrectChoice Os resultados de cada estudo e o efeito combinado
\choice Apenas a metodologia dos estudos
\choice A sequência de passos do estudo
\end{choices}

%-----------------------------------
\question O que indica um fluxograma simplificado de meta-análise?  
\begin{choices}
\choice Os tipos de efeito
\CorrectChoice As etapas do estudo: inclusão, extração de dados, cálculo do efeito e interpretação
\choice O tamanho da amostra
\choice O gráfico de significância
\end{choices}

%-----------------------------------
\question Qual destas afirmações sobre meta-análise é correta?  
\begin{choices}
\choice É útil apenas para estudos qualitativos
\CorrectChoice Permite sintetizar resultados quantitativos e aumenta confiança nas evidências
\choice Sempre substitui a necessidade de análise estatística individual
\choice Ignora heterogeneidade entre estudos
\end{choices}

\end{questions}

\end{document}
